% ============================================================
%  文档类选项（在 main.tex 的 \documentclass[...]{tongjithesis} 中设置）
%  field=science：理工科编号格式 1 / 1.1 / 1.1.1
% ============================================================

% ============================================================
%  封面信息
% ============================================================
\school{电子与信息工程学院}
\major{微电子科学与工程}
\student{2253644}{云发鹏}
\thesistitle{基于电子对抗网络的空地协同部署优化}{}
\thesistitleeng{Air-Ground Cooperative Deployment Optimization Based on Electronic Countermeasure Networks}{}
\thesisadvisor{徐凡}
\thesisdate{2026}{5}{29}

% ============================================================
%  信息说明页
% ============================================================
\infotype{design}

\infoabstract{%
本文研究复杂区域中的空地协同雷达/电子对抗节点部署问题。方法基于多目标粒子群优化-区域分解与坐标变换方法，通过凸分解、垂直交点坐标变换和空地/地地异构传播评估方法，优化节点部署位置，平衡网络覆盖率与干扰强度。仿真结果表明，与原始配置和随机搜索方法相比，所提算法能够生成更具解释性的非支配部署方案，并在边界覆盖和部分质量指标上体现改进，可为复杂区域中的网络覆盖与干扰压制权衡提供参考。
}

\infodrawings{19}
% 毕业设计：主要图纸张数与正文总字数。最终编译后按统计结果更新。
\infowordcount{17463}

\infomaterials{
  \item 论文源文件、参考文献数据库、论文插图和最终PDF文件；
  \item MOPSO-DT算法代码、区域分解与坐标变换代码；
  \item 参数调优、多算法对比、边界实验和消融实验脚本；
  \item JSON/NPZ实验结果文件、实验运行日志和测试用例；
  \item 图表生成脚本、编译日志、结果审计报告和运行说明文档。
}

